% Session 2 of "From Paper to Proof: Formalizing Mathematics in Lean with AI".
% Build: Overleaf with the pdfLaTeX compiler (or pdflatex, twice)
\documentclass[10pt, aspectratio=169]{beamer}
\usepackage{bgsmathlean}
\leanfile{Session2.lean}

\title{BGSMath Graduate Course: From Paper to Proof}
\subtitle{Session 2: From paper to Lean}
\author{Michail Karatarakis}
\institute{Radboud University Nijmegen \& Axiomatic AI}
\date{12 November 2026}

\begin{document}

\maketitle

%============================================================================
\begin{frame}[fragile]{The course}
\sessionplan{2}

\vspace{0.4em}
Each session: a lecture with live demonstrations in Lean. You redo them, and the exercises, at home.
\end{frame}

%----------------------------------------------------------------------------
\begin{frame}[fragile]{Last time, and today}
\begin{keypoint}
Lean checks the \textbf{proof}. You check the \textbf{statement}.
\end{keypoint}

\vspace{0.8em}
\begin{columns}[T]
\begin{column}{0.48\textwidth}
\textbf{\color{leanblue}Last time: reading}
\begin{itemize}
\item Statements written by someone else.
\item Junk values, coercions, quantifier order, precedence.
\item The checklist.
\end{itemize}
\end{column}
\begin{column}{0.48\textwidth}
\textbf{\color{leanblue}Today: writing}
\begin{itemize}
\item Your own definitions and statements.
\item The usual mistakes, and tests that catch them.
\item Letting an AI write the draft, and checking it.
\end{itemize}
\end{column}
\end{columns}
\end{frame}

%============================================================================
\section{Writing definitions}

\begin{frame}[fragile]{A definition, from paper}
\begin{columns}[T]
\begin{column}{0.4\textwidth}
\begin{paper}[Definition]
$f : \mathbb{R} \to \mathbb{R}$ is \emph{bounded} if there is $M$ with $|f(x)| \le M$ for all $x$.
\end{paper}
\end{column}
\begin{column}{0.56\textwidth}
\leanlabel
\lean{bounded_def}
\end{column}
\end{columns}

\vspace{0.6em}
\begin{itemize}
\item \lc{def} gives a name to a statement: \lc{Bounded f} has type \lc{Prop}.
\item The body is the paper sentence, word for word: \emph{there is $M$} becomes \lc{∃ M}, \emph{for all $x$} becomes \lc{∀ x}.
\item The order of the quantifiers is the order of the sentence. $M$ comes first, so it may not depend on $x$.
\end{itemize}
\end{frame}

%----------------------------------------------------------------------------
\begin{frame}[fragile]{Test your definitions\livedemo}
\begin{columns}[T]
\begin{column}{0.55\textwidth}
\lean{bounded_tests}
\end{column}
\begin{column}{0.41\textwidth}
\begin{itemize}
\item A definition is right only if the examples you know come out right.
\item $\sin$ is bounded, $x \mapsto x$ is not. If you cannot prove both, the definition is wrong.
\item One example that \emph{has} the property, one that does not.
\end{itemize}
\end{column}
\end{columns}

\vspace{0.6em}
\begin{keypoint}
\normalsize Write the tests before you use the definition in a theorem.
\end{keypoint}
\end{frame}

%----------------------------------------------------------------------------
\begin{frame}[fragile,t]{Spot the problem (1)}
\begin{paper}[Definition]
A natural number $n$ is \emph{prime} if its only divisors are $1$ and $n$.
\end{paper}

\pause
\begin{columns}[T]
\begin{column}{0.55\textwidth}
\lean{isprime_bad}
\end{column}
\begin{column}{0.41\textwidth}
\begin{itemize}
\item By this definition $1$ is prime. On paper ``$n > 1$'' goes without saying; in Lean it must be written.
\item A test caught it: \lc{IsPrime 1} should have been false.
\end{itemize}
\end{column}
\end{columns}
\end{frame}

%----------------------------------------------------------------------------
\begin{frame}[fragile]{Use Mathlib's definition}
\lean{perfect}

\vspace{0.4em}
\begin{itemize}
\item $n$ is \emph{perfect} if it is the sum of its divisors smaller than $n$: $6 = 1 + 2 + 3$.
\item Why \lc{0 < n}? Without it $0$ is perfect: its set of divisors smaller than $0$ is empty, and the empty sum is $0$.
\item Mathlib's definitions have been reviewed, and their edge cases thought through. If Mathlib has the notion (\lc{Even}, \lc{Nat.Prime}, \lc{Continuous}, \lc{Monotone}, \dots), use it.
\end{itemize}
\end{frame}

%----------------------------------------------------------------------------
\begin{frame}[fragile]{If you write your own, compare}
\lean{mono_agree}

\vspace{0.4em}
\begin{itemize}
\item \lc{Iff.rfl}: the two say the same thing, by definition.
\item When they do not agree by definition, prove the equivalence. The proof is a test of your definition.
\item Then use Mathlib's version in your statements, so that Mathlib's lemmas apply to them.
\end{itemize}
\end{frame}

%============================================================================
\section{Writing statements}

\begin{frame}[fragile]{From paper to statement}
\begin{enumerate}
\item \textbf{Types.} What is each object? A number in $\mathbb{N}$, $\mathbb{Z}$ or $\mathbb{R}$, a function, a set?
\item \textbf{Hypotheses.} Everything the paper assumes, including what it leaves implicit.
\item \textbf{Quantifiers.} What may depend on what?
\item \textbf{The conclusion}, in Mathlib's words: \lc{Set.Icc a b} for $[a,b]$, \lc{ContinuousOn}, \lc{Tendsto}, \dots
\end{enumerate}

\vspace{0.6em}
\begin{keypoint}
\normalsize The same four points as the reading checklist, from the other side.
\end{keypoint}
\end{frame}

%----------------------------------------------------------------------------
\begin{frame}[fragile]{A worked example\livedemo}
\begin{columns}[T]
\begin{column}{0.38\textwidth}
\begin{paper}[Theorem]
A continuous function on $[a, b]$ attains its maximum.
\end{paper}

\begin{itemize}
\item \lc{Set.Icc a b} is $[a, b]$.
\item \lc{ContinuousOn}: continuous on the interval only.
\item \lc{∃ c ∈ s}, \lc{∀ x ∈ s}: quantifiers over a set.
\item Where is \lc{a ≤ b} on paper?
\end{itemize}
\end{column}
\begin{column}{0.58\textwidth}
\leanlabel
\lean{max_icc}
\end{column}
\end{columns}
\end{frame}

%----------------------------------------------------------------------------
\begin{frame}[fragile,t]{Spot the problem (2)}
\lean{max_icc_bad}

\pause
\begin{itemize}
\item If $a > b$, the interval is empty and there is no $c$. Without \lc{hab} the statement is false.
\item On paper the hypothesis was hidden in the notation $[a, b]$.
\end{itemize}

\begin{keypoint}
\normalsize Each hypothesis should be needed: remove it, and the statement should fail.
\end{keypoint}
\end{frame}

%----------------------------------------------------------------------------
\begin{frame}[fragile,t]{Spot the problem (3)}
\begin{paper}[Claim]
For every $x$, \ $x - 1 < x$.
\end{paper}

\pause
\begin{columns}[T]
\begin{column}{0.55\textwidth}
\lean{infer_nat}
\end{column}
\begin{column}{0.41\textwidth}
\begin{goal}
x : ℕ
⊢ x - 1 < x
\end{goal}
\small
\begin{itemize}
\item With no type given, the \lc{1} makes Lean choose $\mathbb{N}$. There the claim fails at $0$.
\item Always write the type. Hover over \lc{x} to check.
\end{itemize}
\end{column}
\end{columns}
\end{frame}

%----------------------------------------------------------------------------
\begin{frame}[fragile]{Test on small cases\livedemo}
\begin{paper}[Claim]
$1 + 2 + \dots + n = \dfrac{n(n+1)}{2}$.
\end{paper}

\lean{sum_eval}

\begin{itemize}
\item \lc{#eval} computes both sides for $n = 0, \dots, 4$. They differ: the statement is wrong.
\item Testing costs nothing, and it works before you have a proof.
\end{itemize}
\end{frame}

%----------------------------------------------------------------------------
\begin{frame}[fragile]{Off by one}
\lean{sum_bug}

\vspace{0.4em}
\begin{itemize}
\item \lc{Finset.range n} is $\{0, \dots, n-1\}$, not $\{1, \dots, n\}$.
\item The second statement is right, and \lc{/} on $\mathbb{N}$ does no harm here: $n(n+1)$ is even. Check that whenever you divide in $\mathbb{N}$.
\end{itemize}
\end{frame}

%----------------------------------------------------------------------------
\begin{frame}[fragile]{Limits, the Mathlib way}
\begin{columns}[T]
\begin{column}{0.4\textwidth}
\begin{paper}[Definition]
$u_n \to a$ if for every $\varepsilon > 0$ there is $N$ with $|u_n - a| < \varepsilon$ for all $n \ge N$.
\end{paper}
\end{column}
\begin{column}{0.56\textwidth}
\leanlabel
\lean{tendsto}
\end{column}
\end{columns}

\vspace{0.4em}
\begin{leanout}
Metric.tendsto_atTop : Tendsto u atTop (𝓝 a) ↔
  ∀ ε > 0, ∃ N, ∀ n ≥ N, dist (u n) a < ε
\end{leanout}

\begin{itemize}
\item Mathlib states limits as \lc{Tendsto u atTop (𝓝 a)}. Write your statements this way, so that Mathlib's lemmas apply.
\item Lemmas like \lc{Metric.tendsto_atTop} translate it back into the $\varepsilon$--$N$ form.
\end{itemize}
\end{frame}

%============================================================================
\section{Checking a statement}

\begin{frame}[fragile]{Can the hypotheses hold?}
\lean{witness}

\vspace{0.4em}
\begin{itemize}
\item Give one example where all the hypotheses hold. If there is none, the theorem says nothing (session 1, \lc{everything_is_42}).
\item This is a test of the statement, so do it before the proof.
\end{itemize}
\end{frame}

%----------------------------------------------------------------------------
\begin{frame}[fragile]{Use it on a case you know}
\lean{special_case}

\vspace{0.4em}
\begin{itemize}
\item Apply the statement to an example where you know the answer.
\item If this is hard, the statement may be stated in a form nobody can use.
\end{itemize}

\vspace{0.6em}
\begin{athome}
\texttt{Exercises2.lean}, parts 1 and 2: write the definitions and statements, then run these tests on them.
\end{athome}
\end{frame}

%----------------------------------------------------------------------------
\begin{frame}[fragile]{Checking a statement: the tests}
\begin{enumerate}
\item \textbf{Read it back.} Translate the Lean into English \emph{without} looking at the paper, then compare.
\item \textbf{Small cases.} \lc{#eval}, \lc{decide}, \lc{norm_num} on examples.
\item \textbf{The hypotheses can hold.} Give an example.
\item \textbf{Each hypothesis is needed.} Remove it, and the statement fails.
\item \textbf{A case you know.} Apply the statement to it.
\item \textbf{The negation.} If you can prove it, the statement is false.
\item \textbf{Mathlib.} If Mathlib has the result, compare with its statement.
\end{enumerate}
\end{frame}

%============================================================================
\section{Drafting statements with AI}

\begin{frame}[fragile]{What an AI can do here}
\begin{columns}[T]
\begin{column}{0.48\textwidth}
\textbf{\color{okgreen}Useful}
\begin{itemize}
\item Finding Mathlib's names: \lc{Set.Icc}, \lc{ContinuousOn}, \lc{Tendsto}.
\item Getting a draft to compile.
\item Translating a Lean statement back into English.
\end{itemize}
\end{column}
\begin{column}{0.48\textwidth}
\textbf{\color{alertred}Check for}
\begin{itemize}
\item The usual mistakes: a missing hypothesis, the wrong type, off by one.
\item A statement changed until it is easy to prove.
\item A definition of its own where Mathlib has one.
\end{itemize}
\end{column}
\end{columns}

\vspace{0.8em}
\begin{keypoint}
\normalsize The AI writes the draft. You are responsible for what it says.
\end{keypoint}
\end{frame}

%----------------------------------------------------------------------------
\begin{frame}[fragile]{A prompt that helps}
\begin{paper}[Prompt]
\small Formalize the following statement in Lean 4 with Mathlib. Use Mathlib's definitions; do not define your own. Do not prove it: end with \texttt{sorry}. Then:
\begin{enumerate}
\item translate your Lean statement back into English;
\item list each hypothesis and say why it is needed;
\item give one example where all the hypotheses hold.
\end{enumerate}
\emph{Statement:} \dots
\end{paper}

\vspace{0.4em}
\begin{itemize}
\item Asking for \lc{sorry} keeps the statement and the proof apart.
\item Steps 1--3 are three of the tests. Check the answers yourself.
\end{itemize}
\end{frame}

%----------------------------------------------------------------------------
\begin{frame}[fragile,t]{Spot the problem (4)}
\lean{draft_circular}

\pause
\begin{leanout}
warning: Variable name `hmono` is not explicitly referenced.
warning: Variable name `hbdd` is not explicitly referenced.
\end{leanout}

\begin{caveat}[Lean accepts this]
The conclusion is one of the hypotheses: \lc{hconv}. The theorem is true and says nothing.
\end{caveat}
\end{frame}

%----------------------------------------------------------------------------
\begin{frame}[fragile]{The fix}
\lean{draft_fixed}

\vspace{0.4em}
\begin{itemize}
\item Without \lc{hconv}, the proof needs a real lemma: \lc{tendsto_atTop_ciSup}, and the limit is $\sup_n a_n$.
\item A hypothesis the proof never uses is a warning sign. Either it is not needed, or the statement is wrong.
\end{itemize}
\end{frame}

%----------------------------------------------------------------------------
\begin{frame}[fragile]{The workflow}
\centering
\begin{tikzpicture}[node distance=0.5cm, every node/.append style={text=ink}]
  \node[layer=paperrule, minimum width=2.9cm] (p) {Paper statement\\[-2pt]{\scriptsize every hypothesis written}};
  \node[layer=leanblue, minimum width=2.9cm, right=0.55cm of p] (d) {AI draft\\[-2pt]{\scriptsize ends with \texttt{sorry}}};
  \node[layer=alertred, minimum width=2.9cm, right=0.55cm of d] (c) {You check\\[-2pt]{\scriptsize checklist and tests}};
  \node[layer=okgreen, minimum width=2.9cm, right=0.55cm of c] (r) {Then the proof\\[-2pt]{\scriptsize session 3}};
  \draw[arr=ink] (p) -- (d);
  \draw[arr=ink] (d) -- (c);
  \draw[arr=ink] (c) -- (r);
  \draw[arr=alertred] (c.south) to[bend left=35] node[lbl=alertred, below] {fix and check again} (d.south);
\end{tikzpicture}

\vspace{1.2em}
\raggedright
\begin{keypoint}
\normalsize Do not ask for a proof until you trust the statement.
\end{keypoint}
\end{frame}

%============================================================================
\section{At home}

\begin{frame}[fragile]{At home}
\begin{columns}[T]
\begin{column}{0.55\textwidth}
\texttt{Exercises2.lean}, in the browser or on your laptop.
\begin{enumerate}
\item Definitions: write the \lc{def}; its tests are below it.
\item Statements: from the paper sentence. The proof is in Mathlib.
\item Find the bug: true, false, or true for a silly reason? Fix the false ones.
\item Your own statement, from session 1. Draft it with an AI, then check it.
\end{enumerate}
\end{column}
\begin{column}{0.41\textwidth}
\begin{itemize}
\item First redo the demos. Their code is in \texttt{Session2.lean}.
\item Part 4: note every change you make to the AI's draft, and why.
\item Send me your statement: in session 4 I formalize some of them live.
\end{itemize}
\end{column}
\end{columns}

\vspace{0.3em}
\begin{keypoint}
\normalsize Today the stating numbers are higher than the proving numbers. That is the point.
\end{keypoint}
\end{frame}

%----------------------------------------------------------------------------
\begin{frame}[fragile]{Next time: AI tools}
\begin{itemize}
\item AI agents that write Lean proofs, in practice.
\item Case studies: how the large formalizations of the last year were done.
\item An agent proving a statement, live. At home, you try it on yours.
\end{itemize}

\vspace{1em}
\begin{keypoint}
\normalsize Before next time: Lean and Mathlib working on your laptop (see the setup sheet), and your statement from part 4 compiling with \lc{sorry}.
\end{keypoint}
\end{frame}

%----------------------------------------------------------------------------
\begin{frame}[fragile]{Resources}
\begin{itemize}
\item J.\ Avigad and P.\ Massot, \emph{Mathematics in Lean}. Its chapters on sets, functions and structures show how Mathlib states things.
\item K.\ Buzzard, \emph{Formalising Mathematics} (Imperial College). Notes on reading and writing statements.
\item The Mathlib documentation, Loogle and LeanSearch: find Mathlib's version before you write your own.
\item The Lean Zulip chat, stream \emph{\#new members}.
\end{itemize}
\end{frame}

\end{document}
